\documentclass[10pt,a4paper]{amsart}
\usepackage[margin=19mm]{geometry}
\usepackage{amsmath,amssymb,mathtools}
\usepackage[scr=rsfs,cal=euler]{mathalfa}
\usepackage{xfrac}
\usepackage[unicode,hidelinks,bookmarks=false]{hyperref}
\newcommand{\ie}{i.e.\ }
\newcommand{\cf}{cf.\ }
\newcommand{\resp}{respectively\ }
\newcommand{\Subsection}{\subsection}
\providecommand{\nobreakdash}{\nobreak\hbox{-}}
\setlength{\emergencystretch}{2em}
\multlinegap0pt
\usepackage{amssymb}
\newtheorem{theorem}{Theorem}
\newcommand{\TT}{\ensuremath{\mathbb{T}}}
\newcommand{\abs}[1]{\ensuremath{\left\lvert{#1} \right\rvert}}
\newcommand{\beq}{\begin{equation}}
\newcommand{\drm}{\ensuremath{\mathrm{d}}}
\newcommand{\eeq}{ \end{equation} }
\newcommand{\htt}{\ensuremath{\hat{t}}}
\newcommand{\no}{\nonumber}
\title{A Multirate Transient Simulation Technique for Circuits described by Integro-Differential Equations}
\author{Hans Georg Brachtendorf, Kai Bittner}
\date{Source: arXiv:2609.25854v1 (CC BY 4.0)}
\begin{document}
\maketitle

\noindent\emph{Source and licence.} A Multirate Transient Simulation Technique for Circuits described by Integro-Differential Equations. Version arXiv:2609.25854v1 (CC BY 4.0), distributed by arXiv under the Creative Commons Attribution 4.0 International licence, http://creativecommons.org/licenses/by/4.0/, as recorded in the arXiv licence metadata for this exact version and shown on the abstract page https://arxiv.org/abs/2609.25854v1. Full licence notice: \texttt{licenses/arxiv-2609.25854-CC-BY-4.0.txt}.\par\smallskip

\noindent\textit{Editorial scope.} Counted unit: the article's theorem theo.4.2 (source0.tex lines 561--571). The article's complete original proof of it is delivered with it (source0.tex lines 573--588, one \texttt{proof} environment of 16 lines). No mathematics is written, repaired or re-derived: the statement and the proof are the article's own lines, in the article's own notation, and no retained line is substituted or re-flowed.  Retained uncounted as the source-local context the proof needs: lines 353--358; lines 360--365; lines 520--530.  Nothing was omitted from any retained span, so the package carries no omission record.  No retained line contains a citation, so the wrapper carries no bibliography and no bibliography entry.  No retained line contains a figure environment, an image-inclusion command or drawing code, so no image is delivered and the package declares no asset dependency.  Editorial formatting only, in addition: an AMS \texttt{amsart} preamble that restates the article's own declarations and macros used by the retained lines, namely the environments \texttt{theorem} on separate counters, together with the standard packages those lines need.  The retained environments print in this document as Theorem 1 in order of appearance, whereas the article prints the same items with the numbering of its own full document.  The complete native source of the article as distributed in the arXiv e-print for this exact version is bundled unchanged as \texttt{sources/211326/source0.tex}.
	\begin{align}
	 & h(t-\tau')\, \hat{x}(\tau',\,\omega(\tau')\,\tau'+\Theta) \no\\
	 &\qquad = \int_\TT
	      \hat{h}(t-\tau',\,\omega(t)\,t+\Theta-t')\,\hat{x}(\tau',\,t')\,\drm t'
	      \label{eq.theo.1}
	\end{align}

	\begin{align}
	 & h(\tau')\,\hat{x}(t-\tau',\,\varphi(t-\tau')+\Theta) = \no\\
	 &\qquad \int_\TT\,
	      \hat{h}(\tau',\,t')\,\hat{x}(t-\tau',\,\omega(t)\,t+\Theta-t')\,\drm t'
	      \label{eq.theo.2}
	\end{align}

\begin{theorem}
    \label{theo.3}
    Let $\omega(\tau)=\omega_0=\text{const.}$
    The equalities \eqref{eq.theo.1} and \eqref{eq.theo.2}
    are fulfilled if
	\beq
	 \hat{h}(\tau,\,\hat{t}) = h(\tau)\,\delta(\hat{t} - \omega_0\,\tau) =
	 \sum_k \frac{1}{2 \pi} h(\tau)\, \mathrm{e}^{-\jmath \, k\, \omega_0 \,\tau}
	 \,\mathrm{e}^{\jmath \,k\, \htt} \label{80}
	\eeq
\end{theorem}

\begin{theorem}
   \label{theo.4.2}
   Let $\omega(\tau)=\omega_0=\text{const.}$  The equality
 	\beq
	 \hat{h}(\tau,\,\hat{t}) = 
	 \sum_k \frac{1}{2 \pi} h(\tau)\, \mathrm{e}^{-\jmath \,  k\, \omega_0 \,\tau}
	 \,\mathrm{e}^{\jmath \,k\, \htt} = h(\tau)\, \delta(\omega_0\,\tau - \htt) \no
%	 \label{80a}
	\eeq
holds.
\end{theorem}

\begin{proof}
 Since
 \[
   \delta(t) = \frac{1}{T}\sum_k \mathrm{e}^{-\jmath \, 2\pi \, k\, t/T}
 \]
 holds on $\TT$, one gets from \eqref{80} the identity
 \begin{align}
 \hat{h}(\tau,\,\hat{t}) &= 
	 \frac{1}{2 \pi} h(\tau)\sum_k  \mathrm{e}^{-\jmath \, k \,\omega_0 \, \tau}
	 \,\mathrm{e}^{\jmath \,k\, \htt} = \frac{1}{\omega_0} h(\tau)\,
	 \delta\left(\tau - \frac{\htt}{\omega_0} \right)	 \no\\
	 &= h(\tau)\sum_k \delta(\omega_0\,\tau - \htt)	 \no
 \end{align}
 where the latter equality results from the identity $\delta(a\,t - t_0) = \frac{1}{\abs{a}} 
 \delta\left(t - \frac{t_0}{a}\right)$.
\end{proof}

\end{document}
